
\documentclass{article}
\usepackage{colortbl}
\usepackage{makecell}
\usepackage{multirow}
\usepackage{supertabular}

\begin{document}

\newcounter{utterance}

\centering \large Interaction Transcript for game `eqbench', experiment `eq\_bench', episode 135 with reflexion{-}llama{-}3.1{-}8b{-}instruct.
\vspace{24pt}

{ \footnotesize  \setcounter{utterance}{1}
\setlength{\tabcolsep}{0pt}
\begin{supertabular}{c@{$\;$}|p{.15\linewidth}@{}p{.15\linewidth}p{.15\linewidth}p{.15\linewidth}p{.15\linewidth}p{.15\linewidth}}
    \# & \multicolumn{2}{c}{Player} && \multicolumn{2}{c}{Game Master} \\
    \hline

    \theutterance \stepcounter{utterance}  
    & & & \multicolumn{4}{p{0.6\linewidth}}{
        \cellcolor[rgb]{0.9,0.9,0.9}{
            \makecell[{{p{\linewidth}}}]{
                \texttt{\tiny{[P1$\langle$GM]}}
                \texttt{You have attempted this same task before and did not succeed. These are the plans you wrote after those attempts.} \\
\texttt{Trial \#1: New plan: Upon re{-}reading the dialogue, I realize that I focused too much on Alessandra's responses and didn't consider Jacques' tone and words. Jacques is not just criticizing Alessandra's cooking, but also her attitude and approach to cooking. He's implying that she's not passionate or dedicated enough to achieve excellence. } \\
\\ 
\texttt{In my next attempt, I will pay closer attention to Jacques' words and tone, and consider how they might affect Alessandra's emotional state. I will also look for more subtle cues in the dialogue, such as body language or tone of voice, that might indicate Alessandra's emotional response.} \\
\\ 
\texttt{I will also consider the context of the conversation. Jacques is a chef, and Alessandra is his employee. This power dynamic might affect Alessandra's emotional response to Jacques' criticism. } \\
\\ 
\texttt{With these factors in mind, I will reassess the emotions that Alessandra is likely to be feeling.} \\
\\ 
\\ 
\\ 
\texttt{Your task is to predict the likely emotional responses of a character in this dialogue:} \\
\\ 
\texttt{Jacques: Alessandra, you've been under{-}seasoning the scallops all night. } \\
\texttt{Alessandra: I've been following the recipe, chef.} \\
\texttt{Jacques: (Laughs) The recipe? That's your problem, Alessandra. You're a slave to the recipe. You need to feel the food, taste it, not just blindly follow instructions.} \\
\texttt{Alessandra: I understand that, chef.} \\
\texttt{Jacques: Do you? Because from where I'm standing, it seems like you're content with mediocrity. } \\
\texttt{Alessandra: That's not fair, Jacques.} \\
\texttt{Jacques: Fair? This isn't about fairness, Alessandra. This is about excellence. } \\
\texttt{Alessandra: I'm doing my best.} \\
\texttt{Jacques: Your best isn't good enough. Not for this kitchen. Not for me.} \\
\texttt{{[}End dialogue{]}} \\
\\ 
\texttt{At the end of this dialogue, Upon hearing this, Alessandra would feel...} \\
\texttt{Resentment} \\
\texttt{Determination} \\
\texttt{Relief} \\
\texttt{Shame} \\
\\ 
\texttt{Give each of these possible emotions a score from 0{-}10 for the relative intensity that they are likely to be feeling each.} \\
\\ 
\texttt{You must output in the following format, including headings (of course, you should give your own scores), with no additional commentary:} \\
\\ 
\texttt{Resentment: <score>} \\
\texttt{Determination: <score>} \\
\texttt{Relief: <score>} \\
\texttt{Shame: <score>} \\
\\ 
\\ 
\texttt{{[}End of answer{]}} \\
\\ 
\texttt{Remember: zero is a valid score, meaning they are likely not feeling that emotion. You must score at least one emotion > 0.} \\
\\ 
\texttt{Your answer:} \\
            }
        }
    }
     \\ \\

    \theutterance \stepcounter{utterance}  
    & \multicolumn{4}{p{0.6\linewidth}}{
        \cellcolor[rgb]{0.561,0.737,0.561}{
            \makecell[{{p{\linewidth}}}]{
                \texttt{\tiny{[P1$\rangle$GM]}}
                \texttt{Resentment: 8} \\
\texttt{Determination: 2} \\
\texttt{Relief: 0} \\
\texttt{Shame: 9} \\
            }
        }
    }
    & & \\ \\

    \theutterance \stepcounter{utterance}  
    & & & \multicolumn{2}{p{0.3\linewidth}}{
        \cellcolor[rgb]{0.9,0.9,0.9}{
            \makecell[{{p{\linewidth}}}]{
                \texttt{\tiny{[GM$|$GM]}}
                \texttt{Resentment: 8} \\
\texttt{Determination: 2} \\
\texttt{Relief: 0} \\
\texttt{Shame: 9} \\
            }
        }
    }
    & & \\ \\

    \theutterance \stepcounter{utterance}  
    & & & \multicolumn{2}{p{0.3\linewidth}}{
        \cellcolor[rgb]{0.9,0.9,0.9}{
            \makecell[{{p{\linewidth}}}]{
                \texttt{\tiny{[GM$|$GM]}}
                \texttt{Resentment: 8} \\
\texttt{Determination: 6} \\
\texttt{Relief: 0} \\
\texttt{Shame: 7} \\
            }
        }
    }
    & & \\ \\

    \theutterance \stepcounter{utterance}  
    & & & \multicolumn{2}{p{0.3\linewidth}}{
        \cellcolor[rgb]{0.9,0.9,0.9}{
            \makecell[{{p{\linewidth}}}]{
                \texttt{\tiny{[GM$|$GM]}}
                \texttt{game\_result = LOSE} \\
            }
        }
    }
    & & \\ \\

    \theutterance \stepcounter{utterance}  
    & & \multicolumn{4}{p{0.6\linewidth}}{
        \cellcolor[rgb]{0.95,0.95,0.95}{
            \makecell[{{p{\linewidth}}}]{
                \texttt{\tiny{[P1$\leftrightarrow$P]}}
                \texttt{Post{-}episode Reflexion (not sent to GM):} \\
\texttt{New plan: Upon re{-}reading the dialogue, I realize that I focused too much on Alessandra's responses and didn't consider Jacques' tone and words. Jacques is not just criticizing Alessandra's cooking, but also her attitude and approach to cooking. He's implying that she's not passionate or dedicated enough to achieve excellence. } \\
\\ 
\texttt{In my next attempt, I will pay closer attention to Jacques' words and tone, and consider how they might affect Alessandra's emotional state. I will also look for more subtle cues in the dialogue, such as body language or tone of voice, that might indicate Alessandra's emotional response.} \\
\\ 
\texttt{However, I also realize that I didn't consider the impact of Jacques' words on Alessandra's self{-}perception. His criticism is not just about her cooking, but also about her character and abilities. This could lead to feelings of inadequacy and self{-}doubt.} \\
\\ 
\texttt{To improve my assessment, I will consider the following:} \\
\\ 
\texttt{1. The severity of Jacques' criticism: He calls Alessandra's cooking "mediocrity" and implies that she's not good enough for his kitchen. This could lead to feelings of shame and inadequacy.} \\
\texttt{2. The tone of Jacques' words: He laughs and uses a condescending tone, which could make Alessandra feel belittled and disrespected.} \\
\texttt{3. The power dynamic: Jacques is a chef and Alessandra is his employee, which could make Alessandra feel vulnerable and defensive.} \\
\\ 
\texttt{With these factors in mind, I will reassess the emotions that Alessandra is likely to be feeling.} \\
\\ 
\texttt{Resentment: 6} \\
\texttt{Determination: 4} \\
\texttt{Relief: 0} \\
\texttt{Shame: 9} \\
            }
        }
    }
    & \\ \\

\end{supertabular}
}

\end{document}
